\section{Sources and contributions}\label{sec:sources}

\paragraph{Sources.} The note \cite{remark}; the Lean source \cite{alpotree} at the archive hash given
in \S\ref{sec:reading}, with Lean~4 \cite{lean4} and Mathlib \cite{mathlib2020}; and the published
literature cited in the text. No other account of this claim was consulted.

\paragraph{Contributions of this article.}
\begin{itemize}
\item The map of the proof from \lean{riemannHypothesis} to the three deep inputs, with each node
  quoted or cited at file:line (\S\S\ref{sec:top}--\ref{sec:connection}, Figure~\ref{fig:top}).
\item The observation that all five exported theorems reduce to Statement~\ref{st:every} by
  Mathlib-level arguments, and an account of how the literature fields are supplied: 195
  declarations take the structure as an argument and 13 reach its fields, and the final construction
  supplies the first and third fields with the same Galois-representation declaration and the middle
  field with \lean{UpperTheory.Zero.LLC.recTExists}.
\item The reading of the analytic half: the prime sum of remark item 5 as the error term of the
  smoothed explicit formula, the Landau-type criterion on windows (Statement~\ref{st:landau}), and the
  reduction of the square-root bound to level-compatible, chain-tied families.
\item The reading of the window bound (Statement~\ref{st:vwin}) as a bound on lengths of
  $\text{Ш}^1(\mathbb Q_\ell/\mathbb Z_\ell(-e))$ over twisted layers of the cyclotomic
  $\mathbb Z_2$-tower, compared with remark items 11 and 17.
\item The comparison of the three deep inputs (Statements~\ref{st:row}, \ref{st:tlift} and the local
  Langlands correspondence) with the published theorems they resemble, and the observation that the
  traced route to Statement~\ref{st:row} uses a theta-lift input and an input supplied by
  \lean{UpperTheory.Gal.Int.Geo.geoOne_holds}, through a weight-one geometric realization.
\item The crosswalk of remark items 1--20 (Table~\ref{tab:remark}), including the finding that the
  Behrend construction lies outside our estimate of the dependency cone.
\item The tooling: a text index of the tree, a reference-graph estimate of the dependency cone with
  controls, and a script that extracts every quoted statement and checks every citation
  (Appendix~\ref{app:method}).
\end{itemize}
